\documentclass{article}

\usepackage[a4paper,margin=1in]{geometry}
\usepackage{zed-csp}

\title{State Schema}

\begin{document}

\maketitle

\section{State Schemas}
We declare variables above the line and constrain their values in the predicate below the line
The schema box looks different than its axiomatic definition.
It has a name so that we can refer to its contents just by naming it.

\begin{schema}{Editor}
    left, right : TEXT
\where
    \#(left \cat right) <= maxsize
\end{schema}

We declare variables above the line and constrain their values in the predicate
below the line. The schema has a name so we can refer to its by its contents.
The schema box is closed to remind do you that the variables declared inside
the box are local. They can only be used in the schema where it was declared.

\end{document}